\chapter{Płytka pośrednicząca w komunikacji}
    Aby usprawnić pobieranie danych przy pomocy kamizelki BRV utworzona została płytka pośrednia, która ma za zadanie odbierać dane przekazywane szeregowo i wysyłać je dalej przy pomocy technologii Bluetooth.
    Udostępniane usługi:
    \begin{itemize}
        \item {\textbf{BRICS BLE} - serwer Bluetooth Low Energy w trybie Peripheral}
        \item {\textbf{BRICS Serial} - wielowątkowa aplikacja pobierająca szergowo dane z kamizeli BRV}
        \item {\textbf{BRICS.service}}
    \end{itemize}
    \\ \\
    Płytka została zrealizowana przy pomocy mikroprocesora Raspberry Pi Zero W, na którym zainstalowany został system operacyjny Raspian 32-bit. Posiada 64 GB pamięci w postaci karty micro SD. Jest zasilana przy pomocy gniazda micro-USB.
    \\ \\
    Obsługą przesyłania danych zajmuje się wielowątkowa aplikacja. Jeden z wątków - BRVSerial - zajmuje się odbieraniem danych na porcie szeregowym i umieszczanie ich w kolejce.
    Natomiast drugi wątek zajmuje się w pierwszej fazie uruchomieniem i zdefiniowaniem serwera Bluetooth LE, a w drugiej fazie pobieraniem danych z kolejki, rozdzielaniem ich na miejsze aby umożliwić ich transport, oraz wysyłaniem pakietów. Wysyłanie pakietów jest realizowane poprzez aktualizacje udostępnianej przez serwer charakterystyki, co powoduje wyemitowanie powiadomienia do wszystkich podłączonych stacji.
\section {Struktura serwera Bluetooth}
    Serwer Bluetooth LE Peripheral zgodnie ze standardem GATT składa się z:
    \begin{itemize}
        \item {\textbf{Serwisu}}
        \item {\textbf{Charakterystyki} - wartości emitowanej przrez serwe, która może ulec zmianie - w tym wypadku obecny pakiet danych}
    \end{itemize}
    
